% Chapter 12 draft: persistent sequence indexes.
% Research chronology and mandatory DOC0287--DOC0312 audit:
%   research_audit/ch12_docs.md
% Current-source anchors:
%   backend/sequences/indexing/sequence_index_type.cppm
%   backend/sequences/indexing/{minimizer_hashmap_index,
%       minimizer_sequence_index,minimizer_hashmap_format6,mpf5_image_v2}.cppm
%   backend/pg/storage_pg/indexing/{index_admin_pg,index_chunked_storage_pg,
%       index_residency_pg,sorted_pair_runs_pg}.cppm
%   backend/pg/scheduler/dispatch_processor.cppm

\section{An index is derived database state}
\label{sec:index-storage-role}

A standalone mapper can treat its index as a file generated from another file.
The user decides when to rebuild it, gives the resulting pathname to the next
process, and accepts that the mapper's command-line parameters are part of the
file's meaning. Inside PostgreSQL this model is insufficient. A BioSeqDB index
belongs to a particular physical relation and sequence column. It was produced
from one source epoch, alphabet, minimizer scheme, value width, \(k\), \(w\),
occurrence threshold, and backend ABI. A transaction may replace it while
another transaction is still reading the previous version. A failed build must
not leave an apparently complete object. A query that requests a strictly
compatible persisted index must not silently obtain a recomputed approximation.

For these reasons BioSeqDB treats a persistent sequence index as
\emph{derived database state with an explicit identity}. The data plane may be
a PostgreSQL relation or, for the newer MPF path, an immutable mapped file, but
the control plane remains inside PostgreSQL. Catalog state decides which
generation is current; a snapshot determines which descriptor a reader can
see; compatibility metadata decides whether the generation can be consumed;
and memory permits bound the pieces admitted into a backend. This is more
machinery than a conventional \texttt{.mmi} file, but it solves a different
problem: making indexed alignment participate safely in transactional SQL.

The index subsystem is deliberately separate from the base sequence storage
described in Chapter~11. The base relation owns rows and sequence locators.
An index sidecar owns a reproducible transformation of those sequence bytes.
Deleting or replacing the index does not change the source sequence; changing
the source epoch makes an old index ineligible. It is also separate from
\texttt{bioseq\_cache}. The cache remembers a completed SQL request under an
execution identity, whereas a minimizer index remembers a reusable search data
structure under a source and algorithm identity.

\begin{figure}[ht]
  \centering
  \begin{tikzpicture}[
      node distance=0.75cm and 0.9cm,
      box/.style={draw,rounded corners,align=center,minimum height=.82cm,
                  minimum width=2.5cm,fill=gray!7},
      strong/.style={box,fill=blue!10,thick},
      >=Latex,font=\small]
    \node[strong] (base) {BioSeq table\\sequence column};
    \node[box,right=of base] (identity) {physical relation identity\\column + source epoch};
    \node[box,right=of identity] (policy) {index policy\\name + update mode};
    \node[strong,below=1.05cm of identity] (head) {current-generation\\descriptor};
    \node[box,below left=.95cm and .45cm of head] (pg) {
      PostgreSQL relation\\provider \(=0\), default};
    \node[box,below right=.95cm and .45cm of head] (mpf) {
      immutable mapped file\\provider \(=1\), opt-in MPF};
    \node[box,below=1.0cm of pg] (select) {selective reader\\and residency cache};
    \node[box,below=1.0cm of mpf] (accel) {CPU or accelerator\\runtime view};

    \draw[->,thick] (base)--(identity);
    \draw[->,thick] (identity)--(policy);
    \draw[->,thick] (policy.south) |- (head.east);
    \draw[->,thick] (identity.south)--(head.north);
    \draw[->,thick] (head)--(pg);
    \draw[->,thick] (head)--(mpf);
    \draw[->,thick] (pg)--(select);
    \draw[->,thick] (mpf)--(accel);
    \draw[->,thick] (select)--(accel);
  \end{tikzpicture}
  \caption{Persistent-index control and data planes. The descriptor visible
  through PostgreSQL selects a provider and binds the payload to source and
  algorithm identities; the current provider setting is never consulted to
  reinterpret an existing generation.}
  \label{fig:index-control-data-plane}
\end{figure}

\subsection{The compile-time registry and SQL policy}
\label{sec:index-registry}

The public abstraction begins with
\texttt{SequenceIndexType}. It describes physical granularity rather than a
specific algorithm:

\begin{enumerate}
  \item \texttt{CHUNKED\_PER\_FRAGMENT} associates one or more chunks with a
  source fragment and an application-defined sub-key;
  \item \texttt{INLINE\_PER\_SEQUENCE} stores one value for a sequence; and
  \item \texttt{CHUNKED\_PER\_TABLE} stores chunks belonging to a whole table.
\end{enumerate}

Index traits bind an algorithm to one of these layouts at compile time. The
PostgreSQL registry currently includes
\texttt{MinimizersPerFragmentIndex} (MPF),
\texttt{MinimizerSequenceMergedIndex} (MPS), and
\texttt{SymbolCountsIndex}; a table-level count index exists only in test
builds. Both minimizer implementations use the fragment-chunk storage adapter,
while symbol counts use the per-sequence adapter. Static assertions verify
that each registered algorithm and storage class agree about their layout.
This prevents a name added to SQL from accidentally selecting a class whose
row key has different semantics.

\begin{table}[ht]
  \centering
  \small
  \caption{Current registered indexes and default policy. ``Automatic'' means
  that the policy is enabled on insertion; it does not mean that every read is
  forced to use the resulting representation.}
  \label{tab:index-registry}
  \begin{tabularx}{\linewidth}{p{.22\linewidth}p{.20\linewidth}
                               p{.20\linewidth}X}
    \toprule
    \textbf{Index name} & \textbf{Layout} & \textbf{Default update} &
    \textbf{Role}\\
    \midrule
    \texttt{minimizers\_per\_fragment} &
    chunked per fragment & on insert &
    Default active read index; preserves fragment-native construction and
    permits direct assembly of the target lookup.\\
    \texttt{minimizers\_per\_sequence} &
    chunked per fragment & manual &
    Merged sequence-wide, hash-partitioned buckets with a directory for
    query-selective loading.\\
    \texttt{symbol\_counts} &
    inline per sequence & on insert &
    Compact metadata used by sequence-aware planning and filtering rather than
    anchor lookup.\\
    \bottomrule
  \end{tabularx}
\end{table}

Update policy and read policy are intentionally distinct. An index may be
\texttt{UPDATE\_ON\_INSERT} or \texttt{UPDATE\_MANUALLY}. At read time the
query contract can allow a persisted index, require its absence, or require a
compatible persisted index. The default active minimizer read name is MPF.
The setting that permits stored-index reads controls whether persistence is
consulted; it does not dictate the transient in-memory shape used after
loading. This separation permits experiments with MPS or accelerator views
without changing the SQL-visible index name, and it supplies a strict mode for
experiments that are intended to measure cold persistence rather than hidden
recomputation.

\section{Two logical minimizer organizations}
\label{sec:index-logical-layouts}

\subsection{Minimizers per fragment}
\label{sec:index-mpf}

MPF is closest to the way sequence data enters the pipeline. Each source
fragment is sketched independently, with the overlap rules described in
Chapter~5, and serialized as a
\texttt{MinimizerHashmapFragmentIndex}. Its in-memory representation has two
contiguous components:

\begin{itemize}
  \item a hash table from minimizer hash to a packed
  \((\mathit{offset},\mathit{count})\) pair; and
  \item a positions array in which an occurrence is encoded as
  \((\mathit{globalPosition}\ll1)\mathbin{\vert}\mathit{strand}\).
\end{itemize}

The version-2 legacy serialization begins with a 32-byte header whose magic is
\texttt{0xB1000001}. Every occupied logical bucket is a 16-byte record
\(\{\texttt{u64 hash},\texttt{u32 offset},\texttt{u32 count}\}\), followed by
the occurrence vector. This representation is attractive at build time:
fragments can be produced without one global sort, and each unit has an
independent source key and checksum boundary. At query time, however, multiple
fragments must be assembled or adapted into the lookup shape expected by the
pipeline. The current fragment-native path reduces that conversion and can
derive either a CPU lookup or a GPU-resident view without changing the source
sequence.

\subsection{Minimizers per sequence}
\label{sec:index-mps}

MPS was introduced to remove the cost of repeatedly merging all fragment
indexes for a long logical sequence. Instead of one serialized hash table per
fragment, it stores hash-ordered buckets for the whole sequence. A compact
directory maps hash intervals to storage sub-keys:
\[
  D_i=(h_i^{\mathrm{lo}},h_i^{\mathrm{hi}},s_i),
\]
where each current version-2 entry occupies 24 bytes
\(\{\texttt{u64 lo},\texttt{u64 hi},\texttt{i64 sub\_key}\}\). The current
directory header is 40 bytes and uses magic \texttt{0xB1000002}. Directory and
bucket sub-keys occupy a reserved namespace so that temporary runs, position
extents, dense seals, generation heads, and the final directory cannot collide.

MPS construction scans 32,768-base sequence ranges with the required halo,
extracts globally positioned minimizer pairs, sorts them into runs, and
performs a \(k\)-way merge. Buckets target approximately 8 MiB. An unusually
frequent hash is allowed to cross a physical chunk boundary so that one heavy
hitter cannot force unbounded memory. Most importantly, the directory is
written last. Its presence is a completeness witness: readers do not infer a
valid MPS index from a prefix of buckets left by a failed build.

The directory makes the cold query path selective. The reader first
deduplicates the query's minimizer hashes. Binary search over the directory
identifies only the intersecting bucket sub-keys. Those buckets are loaded,
global occurrence limits are applied, positions are clipped to the requested
target interval, and a transient lookup is assembled. In a DOC0301 experiment
the query selected 4,227,179 of 523,428,324 target minimizers, or
0.8078 percent. That result demonstrated the value of query selection, but the
specific A16 path was transient and persisted nothing; it is not evidence for
a durable MPS build.

\begin{figure}[ht]
  \centering
  \begin{tikzpicture}[x=1cm,y=1cm,>=Latex,font=\small]
    \node[draw,rounded corners,fill=blue!9,minimum width=3.0cm,
          minimum height=.8cm] (q) at (1.7,4.6)
          {distinct query hashes};
    \node[draw,rounded corners,fill=orange!12,minimum width=3.0cm,
          minimum height=.8cm] (d) at (6.0,4.6)
          {authenticated directory};
    \draw[->,thick] (q)--node[above]{binary search}(d);

    \draw[fill=gray!14] (2.10,3.10) rectangle (3.55,3.82);
    \node[font=\scriptsize] at (2.825,3.46) {\([0,20)\)};
    \draw[fill=green!18] (3.65,3.10) rectangle (5.10,3.82);
    \node[font=\scriptsize] at (4.375,3.46) {\([20,40)\)};
    \draw[fill=gray!14] (5.20,3.10) rectangle (6.65,3.82);
    \node[font=\scriptsize] at (5.925,3.46) {\([40,60)\)};
    \draw[fill=green!18] (6.75,3.10) rectangle (8.20,3.82);
    \node[font=\scriptsize] at (7.475,3.46) {\([60,80)\)};
    \draw[fill=gray!14] (8.30,3.10) rectangle (9.75,3.82);
    \node[font=\scriptsize] at (9.025,3.46) {\([80,100)\)};
    \node[anchor=east] at (1.8,3.46) {hash buckets};

    \draw[->,thick,green!45!black] (4.38,3.10)--(4.38,2.45);
    \draw[->,thick,green!45!black] (7.48,3.10)--(7.48,2.45);
    \node[draw,rounded corners,fill=green!12,align=center,
          minimum width=5.1cm,minimum height=.85cm] (selected) at (5.93,1.98)
          {selected controls and positions only\\global occurrence filter};
    \draw[->,thick] (selected)--(5.93,.95);
    \node[draw,rounded corners,fill=purple!9,align=center,
          minimum width=4.3cm,minimum height=.85cm] at (5.93,.50)
          {query-specific transient lookup\\with clipped target positions};
  \end{tikzpicture}
  \caption{Query-selective MPS loading. The directory is small enough to
  authenticate and search as a whole; only buckets intersecting the query hash
  set are admitted. Format~6 generalizes the same principle to separate
  directory, control, and postings streams.}
  \label{fig:mps-selective-load}
\end{figure}

\subsection{Why both organizations remain useful}

MPF and MPS are not a simple old/new pair. MPF aligns with the fragment-native
storage and ingestion pipeline, avoids a sequence-wide external sort, and is
the current default. MPS moves more work to construction so a query can avoid
touching unrelated hash ranges. The appropriate choice depends on source
fragmentation, query selectivity, update frequency, and whether the dominant
cost is construction, cold admission, or runtime conversion. This is why the
registry keeps both names and why persistence preference is explicit: if both
deferred-persistence policies are enabled, current dispatch chooses MPF first
rather than racing two incompatible writes.

\section{Bounded construction and publication}
\label{sec:index-construction}

\subsection{The failure of an apparently simple global sort}

The first sequence-wide builder accumulated a large vector of
\((\mathit{hash},\mathit{position})\) pairs, sorted it, and emitted buckets.
That algorithm is easy to state and performs well while the vector fits in
memory. It is not a safe PostgreSQL operator for multi-gigabase targets. The
early 4 Gbp external-sort experiments recorded approximately 80 GB of
temporary storage for dna4 and 93 GB for dna5; a 6 Gbp run exceeded 120 GB.
The lesson was not merely to increase a buffer. A database operator needs a
resource envelope that does not depend on the target being friendly.

DOC0299 replaced the unbounded phase with sorted runs backed by PostgreSQL
\texttt{BufFile}. The wire pair is exactly 16 bytes:
\[
  P=(\texttt{u64 kmer},\texttt{u64 encodedPositionAndStrand}).
\]
Pairs first occupy a memory-bounded run. A full run is sorted and may remain
resident only while the run-residency budget permits; otherwise it is spilled.
The merge opens a bounded number of readers and combines runs in passes. The
current sequence builder also spills the accumulator for a single extremely
frequent \(k\)-mer, so neither the number of runs nor a heavy hitter can bypass
the memory contract.

\begin{algorithm}[ht]
  \caption{Bounded MPS construction and directory-last publication}
  \label{alg:mps-build}
  \small
  \begin{algorithmic}[1]
    \Require source sequence identity \(I\), parameters \(k,w\), memory budget
    \State reserve a new build identity and snapshot the selected provider
    \For{each 32,768-base source range with minimizer halo}
      \State sketch the range and convert local occurrences to global positions
      \State append 16-byte \((hash,encodedPosition)\) pairs to the active run
      \If{the run budget is full}
        \State sort the run; retain it only if a residency permit is available
        \State otherwise spill it to a PostgreSQL \texttt{BufFile}
      \EndIf
    \EndFor
    \State finish and sort the last run
    \While{more runs exist than can be opened simultaneously}
      \State merge one bounded fan-in group into a new run
    \EndWhile
    \State merge remaining runs in hash order
    \State form approximately 8 MiB buckets; spill a heavy-hash accumulator
          when necessary
    \State write every bucket with length, checksum, source, and compatibility
          identity
    \State write the ordered directory \textbf{last}
    \State validate the complete unpublished generation
    \State atomically publish its descriptor as the current generation
  \end{algorithmic}
\end{algorithm}

Bounded does not automatically mean fast. DOC0300's first implementation used
15 sorted runs, of which 14 were resident and one was file-backed. It retained
3.14 GiB and spilled 4.976 GiB, yet the measured region took 455.452 s
(456.715 s wall time). PostgreSQL accounted 33.148 GiB of temporary traffic,
440,117 temporary files, and six merge passes. This result was valuable
precisely because it separated semantic success from a viable build: the
algorithm completed under a bound, but its file granularity and merge
amplification were unacceptable.

\subsection{Synchronous and identity-only asynchronous work}

An index build performs two qualitatively different kinds of work. Sketching,
sorting, and serializing are ordinary C++ computation. Opening relations,
registering snapshots, reading tuples, and publishing catalog state are
PostgreSQL operations that must run in a valid backend or background worker.
Earlier scheduler proposals risked obscuring this boundary by treating the
build as a generic task payload.

The current asynchronous MPS route uses an \emph{identity-only job}. The
submitting backend transfers stable identifiers and parameters, not live
PostgreSQL pointers, relation handles, snapshots, or memory-context objects.
The PostgreSQL worker rechecks relation identity and policy, registers its own
snapshot, reopens the source, and rebuilds the MPS representation. The mapped
provider is rejected for this MPS asynchronous route. Current settings also
separate \texttt{index\_async\_persist} from
\texttt{recompute\_async\_persist}; both must permit the deferred write, and
the production gate is restricted to the supported \(k=15,w=10\) profile.

This model has an intentional cost: the worker reads and sketches the source
again instead of consuming a large staged payload. In exchange, cancellation,
transaction abort, and backend exit cannot leave a detached task holding
backend-local state. The design follows a general rule used throughout
BioSeqDB: transfer immutable values across the scheduler boundary and reacquire
PostgreSQL capabilities on the thread that owns them.

\section{Physical persistence in PostgreSQL}
\label{sec:index-physical-persistence}

\subsection{Legacy sidecar relations}
\label{sec:index-legacy-sidecar}

The original storage adapters materialize ordinary PostgreSQL relations in the
\texttt{bioseq\_internal} schema. Their names are derived from the physical
identity of the parent relation together with the index name, bucket, and shard.
This is more robust than deriving a name only from the user-visible table name:
rename does not change the relation identity, and drop/recreate does not make a
new table inherit the previous table's index accidentally.

The three generic layouts have deliberately small schemas:

\begin{table}[ht]
  \centering
  \scriptsize
  \caption{Legacy physical adapters. The minimizer implementations use the
  first row; the remaining rows show why the abstraction is not
  minimizer-specific.}
  \label{tab:index-legacy-schemas}
  \begin{tabularx}{\linewidth}{p{.18\linewidth}p{.47\linewidth}X}
    \toprule
    \textbf{Layout} & \textbf{Columns and primary key} &
    \textbf{Storage properties}\\
    \midrule
    Chunked per fragment &
    \texttt{query\_key BIGINT},
    \texttt{sub\_key BIGINT}, \texttt{chunk\_no INT},
    \texttt{chunk\_bytes INT}, \texttt{blob BYTEA};
    primary key \((query\_key,sub\_key,chunk\_no)\) &
    Heap relation; \texttt{blob} uses external storage. Supports MPF and MPS
    directories/buckets.\\
    Inline per sequence &
    \texttt{query\_key BIGINT}, \texttt{chunk\_bytes INT},
    \texttt{blob BYTEA}; primary key \((query\_key)\) &
    One logical object per sequence, with external storage for the byte value.\\
    Chunked per table &
    \texttt{chunk\_no INT}, \texttt{chunk\_bytes INT},
    \texttt{blob BYTEA}; primary key \((chunk\_no)\) &
    Whole-table abstraction retained and tested, but no production index
    currently uses it.\\
    \bottomrule
  \end{tabularx}
\end{table}

Chunking has two purposes. First, a serialization can exceed PostgreSQL's
practical datum limits. Second, a reader can allocate, checksum, and release
bounded units. The legacy maximum is one GiB minus a 64 MiB safety margin, but
newer generation formats deliberately choose much smaller physical payloads.
Lengths in the row are checked against the actual bytea payload before a chunk
is accepted. The application-level sub-key is not a physical tuple number:
for MPS it identifies directory or hash-range material, while MPF can use it to
identify logical sections of a fragment image.

The first update path used replace-in-place semantics. It cleared an old
sidecar and wrote the replacement. Directory-last publication prevented a
partial MPS build from looking complete, but readers and writers still shared
one mutable physical namespace. This model remains useful for the generic
legacy adapters; it is not the strongest lifecycle now available.

\subsection{Immutable generations}
\label{sec:index-generations}

The generation design changes the unit of publication. A writer allocates a
positive generation identifier and appends all of its rows without modifying
the generation visible to current readers. It validates the finished object
and only then atomically replaces a distinguished generation-zero head. The
old descriptor becomes a predecessor or retired generation rather than being
overwritten under a reader.

The PostgreSQL-relation provider uses an eight-column row:
\[
\begin{split}
 (&\texttt{query\_key},\texttt{generation\_id},
   \texttt{position\_segment},\texttt{sub\_key},\\
  &\texttt{chunk\_no},\texttt{chunk\_bytes},
   \texttt{checksum},\texttt{blob}).
\end{split}
\]
The relation is \texttt{UNLOGGED}, has heap fill factor 100, and stores
\texttt{blob} as \texttt{PLAIN}. Its single unlogged B-tree primary index has
key
\[
  (\texttt{generation\_id},\texttt{query\_key},\texttt{sub\_key},
   \texttt{position\_segment},\texttt{chunk\_no})
\]
and fill factor 90. Creation uses the source relation's tablespace. The reader
inspects the relation descriptor and rejects a sidecar whose number, order,
types, nullability, storage modes, access method, or primary-key definition do
not match the expected schema.

The choice of an unlogged relation is deliberate. A minimizer index is
reproducible from durable sequence data, and logging every derived payload byte
would duplicate write amplification. Unlogged is not the same as unchecked:
generation identity, manifest data, lengths, and checksums still detect partial
or misaddressed content. After an unlogged reset, the generation must be
rebuilt rather than treated as valid. The \(8064\)-byte Format-5 row payload
was selected so a PLAIN bytea and its tuple header fit on one ordinary
PostgreSQL page, avoiding a TOAST fetch for every small cold-read unit.

\begin{figure}[ht]
  \centering
  \begin{tikzpicture}[
      node distance=.65cm and .95cm,
      box/.style={draw,rounded corners,align=center,minimum height=.75cm,
                  minimum width=2.55cm},
      >=Latex,font=\small]
    \node[box,fill=gray!12] (old) {generation \(g\)\\currently visible};
    \node[box,fill=orange!12,right=2.1cm of old] (new) {
      generation \(g+1\)\\unpublished};
    \node[box,fill=blue!12,below=1.0cm of new] (validate) {
      validate all rows, manifest,\\checksums and compatibility};
    \node[box,fill=green!12,below=1.0cm of old] (head) {
      generation-zero head\\atomic descriptor replace};
    \node[box,fill=purple!10,below=1.0cm of head] (retire) {
      predecessor retained\\then retired after visibility horizon};

    \draw[->,thick] (new)--(validate);
    \draw[->,thick] (validate)--node[above,sloped]{publish only on success}(head);
    \draw[->,thick] (old)--(head);
    \draw[->,thick] (head)--(retire);
    \draw[->,dashed,thick] (retire.east) -| node[pos=.25,right]
      {bounded reaper} (new.south);
  \end{tikzpicture}
  \caption{Copy-on-publish generation lifecycle. A failed build can leave
  removable unpublished material, but it cannot make a partial generation
  current. Readers that resolved the predecessor under an older snapshot may
  finish before reclamation.}
  \label{fig:index-generation-lifecycle}
\end{figure}

\subsection{The optional mapped-file provider}
\label{sec:index-mapped-provider}

The storage history contains an important reversal. DOC0381 rejected an mmap
fallback under a strict interpretation that every index byte had to reside in
a PostgreSQL relation. Its diagnosis was nevertheless decisive: a cold
\texttt{just\_index} run spent 103.839560417 s in index-and-map work versus
17.168589792 s for minimap2, a 6.048x region disadvantage. Of the BioSeqDB
time, 62.519740125 s read 4,463,431,732 payload bytes from 562,075 rows and
41.319820292 s remained outside raw payload fetch. The old accounting loop
made 103,003 passes over the extent set---about 1.768 billion visits---and a
second checksum pass reread the same 4.46 GB. The primary-index route then
performed roughly 562,000 heap fetches. PostgreSQL ownership alone did not
make this tuple shape an efficient large immutable image.

DOC0383 reframed ownership more precisely. PostgreSQL 17 does not expose a
supported extension hook for a custom storage manager, so forcing a
tuple-per-page emulation was not the only safe option. BioSeqDB could keep the
\emph{lifecycle and authority} in PostgreSQL while placing the immutable MPF
data plane in an extension-owned file. DOC0384 records that implementation.
\texttt{postgres\_relation} remains provider value 0 and the default;
\texttt{mapped\_file} is provider value 1 and opt-in. A generation descriptor
stores its provider, and every read obeys that stored value rather than the
current configuration setting. Changing a GUC can affect a future build but
cannot reinterpret an existing descriptor.

The mapped MPF image has an explicit little-endian MPF5 version-2 layout and a
4096-byte header. Unlike a dump of native C++ objects, its offsets, lengths,
alignment, reserved bytes, and integer widths are format contracts. The open
path validates the header and identity without performing an eager CRC over
several gigabytes. When a selected extent is used, its address-bound checksum
is verified. Thus a cold reader pays integrity cost for the material it
actually admits rather than hashing the entire index twice before the first
lookup.

Publication follows a filesystem protocol designed to make the descriptor
truthful:

\begin{enumerate}
  \item create a private temporary file with mode 0600 in the extension-owned
  database/tablespace directory;
  \item write the complete image, finalize its header, and synchronize it;
  \item publish with a no-replace operation so an existing generation cannot
  be silently overwritten;
  \item synchronize the containing directory, record device/inode identity,
  and make the immutable artifact mode 0400; and
  \item publish the PostgreSQL descriptor transactionally.
\end{enumerate}

Mappings are registered with a PostgreSQL \texttt{ResourceOwner}, so error
unwinding closes them even when normal C++ control flow is bypassed. Transaction
callbacks use a fixed-capacity action queue rather than allocating or throwing
inside a callback; prepared transactions are rejected while such actions are
armed. Replacement records a full transaction identifier. The reaper waits
until PostgreSQL's global visibility test says that identifier is removable,
then reclaims the retired descriptor and artifact. Its work is bounded per
pass: at most 64 descriptors, 4096 relation rows, 64 MiB of relation payload,
64 files, and 64 orphan candidates. Readers remain lock-free after resolving
the descriptor under one registered snapshot; writers, replacement, clear,
and reaping are serialized by an advisory transaction lock.

This design should not be described as “an arbitrary external cache file.”
The file is immutable, privately named, identity checked, transactionally
described, visibility-delayed on deletion, and unable to select itself. It is
an extension-owned data plane governed by PostgreSQL. Conversely, it should
not be described as PostgreSQL heap storage: backups, replication, crash
handling, and administration must account for the provider distinction.

\section{Formats 5 and 6: from complete images to selective streams}
\label{sec:index-formats}

\subsection{Shape-aware Format 5}

The generation codec evolved through several shapes. Format~3 used sorted
dense controls with global 64-bit positions. Format~4 grouped 32-bit hashes
and local 32-bit positions. Format~5 stores a grouped, structure-of-arrays
shape suitable for direct CPU or accelerator assembly. Its preferred logical
segment is approximately 8 MiB, while the PostgreSQL relation provider splits
that segment into 8064-byte physical rows. A 64-byte generation head points to
manifests; Format-5 manifests and directory entries carry the additional
section sizes and shape metadata required to authenticate the image.

The essential change is that the persistent shape is no longer merely “some
bytes that deserialize into a C++ hash map.” It is a specified arrangement of
controls and positions that can be range checked, checksummed, and adapted
without rebuilding the original container. The same logical Format-5
generation can be backed by PostgreSQL rows; the MPF mapped image carries a
closely related direct layout with a provider-specific outer header.

\subsection{Hash-major Format 6}
\label{sec:index-format6}

Format~6 makes query selectivity a physical property. It separates three
storage-neutral streams:

\begin{enumerate}
  \item an authenticated directory over hash ranges;
  \item control blocks containing up to 4096 hashes and global occurrence
  counts; and
  \item postings extents containing positions for hashes that survived the
  occurrence policy.
\end{enumerate}

The fixed head is 224 bytes and the manifest is 336 bytes; a directory entry
is 80 bytes. PostgreSQL payload rows remain 8064 bytes and the reader coalesces
authorized ranges into read slabs of at most 8 MiB. It first validates the
complete directory, intersects query hashes with directory ranges, reads only
the corresponding control blocks, rejects hashes above
\texttt{max\_occ}, and authorizes only the resulting postings intervals.
Exact-count exception data preserves correctness where a compact count field
would otherwise saturate.

\begin{figure}[ht]
  \centering
  \begin{tikzpicture}[x=1cm,y=1cm,>=Latex,font=\scriptsize]
    \draw[rounded corners,thick,fill=blue!9] (0,3.9) rectangle (12.6,4.8);
    \node[anchor=west,font=\bfseries] at (.2,4.55) {Format 6 head and manifest};
    \node[anchor=west] at (.2,4.18) {
      source epoch, generation, scheme, stream lengths, checksums,
      compatibility digest};

    \draw[rounded corners,thick,fill=orange!12] (0,2.65) rectangle (12.6,3.55);
    \node[anchor=west,font=\bfseries] at (.2,3.30) {Directory};
    \node[anchor=west] at (.2,2.93) {
      hash range \(\rightarrow\) control block and postings authorization};

    \foreach \x/\c in {0/gray!15,1/green!18,2/gray!15,3/green!18,4/gray!15} {
      \draw[fill=\c] (.2+\x*2.45,1.45) rectangle (2.45+\x*2.45,2.25);
      \node at (1.325+\x*2.45,1.85) {control \x};
    }
    \node[anchor=east,font=\bfseries] at (-.05,1.85) {controls};

    \foreach \x/\w/\c in
      {0/1.6/gray!15,1.6/3.0/green!18,4.6/1.0/gray!15,
       5.6/4.1/green!18,9.7/2.55/gray!15} {
      \draw[fill=\c] (.2+\x,.2) rectangle (.2+\x+\w,1.0);
    }
    \node[anchor=east,font=\bfseries] at (-.05,.60) {postings};

    \draw[->,thick,green!45!black] (3.875,1.45)--(3.875,1.03);
    \draw[->,thick,green!45!black] (8.775,1.45)--(8.775,1.03);
  \end{tikzpicture}
  \caption{Format~6 separates navigation, hash controls, and variable-length
  postings. Green regions illustrate the blocks authorized by one query;
  unselected postings need not enter the backend.}
  \label{fig:format6-streams}
\end{figure}

The current implementation routes the supported MPS dna4 \(k=15,w=10\)
profile to this format only when its compatibility registration and grouped
generation path are enabled. MPF remains on Format~5. The grouped persisted
generation setting defaults to false in the current source, so Format~6 is an
implemented specialized route, not the unconditional default. Unsupported or
unregistered profiles fall back to the compatible Format-5 path rather than
guessing a wire identity.

\subsection{A rejected per-query occurrence merge}

Selective loading introduces another temptation: compute target occurrence
counts separately for every query and merge them on demand. DOC0385 compared
that candidate with the compact global query-occurrence radix used by the
current path. On its canonical case
\((N,R,U)=(7{,}440{,}026,7{,}258{,}660,7{,}258{,}562)\), the compact global
method took 267,837,208 ns and 206,618,624 bytes peak RSS, while the per-query
merge took 516,441,416 ns and 294,830,080 bytes. The candidate was slower and
larger on every measured workload and was rejected. This is an example of a
local-looking optimization that duplicates work across the SQL batch; the
global radix amortizes it while retaining exact occurrence policy.

\section{Compatibility is part of correctness}
\label{sec:index-compatibility}

Two index images can have identical SQL schemas and still be mutually
unreadable. A 32-bit position image cannot represent a target that required
64-bit positions. A hash produced by the ntHash SIMD scheme does not match the
legacy minimap-style scheme. A GPU kernel compiled for one control layout
cannot safely interpret another merely because both contain \(k=15\)
minimizers. BioSeqDB therefore authenticates meaning as well as bytes.

The compatibility contract pins at least:

\begin{itemize}
  \item alphabet identifier and storage width;
  \item minimizer \(k\) and window \(w\);
  \item scheme identifier, scheme version, and producer backend;
  \item position width and stored shape/format;
  \item occurrence-limit semantics; and
  \item accelerator ABI when the persistent shape is consumed directly by a
  device path.
\end{itemize}

These fields contribute to a compatibility digest carried by the generation
metadata. The reader also binds the generation to the expected relation,
column/query key, source epoch, and generation identifier. It checks manifest
ordering, section lengths, address arithmetic, per-section or per-extent CRCs,
and the content identity appropriate to the format. Reserved bytes must have
their specified value. Integer addition and multiplication are checked before
forming a span, so a corrupt count cannot wrap into an apparently valid
address.

The consequence for SQL semantics is deliberate. A read policy of
\texttt{require\_compatible\_persisted} fails closed on absence or mismatch.
It may not rebuild an index and continue, because that would invalidate a cold
load benchmark and can violate a caller's operational requirement. An
\texttt{allow} policy may fall through to recomputation. A
\texttt{require\_absent} policy supplies the converse experimental contract:
the query must exercise the no-index route.

\begin{table}[ht]
  \centering
  \small
  \caption{Examples of mismatches and the safe response.}
  \label{tab:index-mismatch-response}
  \begin{tabularx}{\linewidth}{p{.28\linewidth}p{.30\linewidth}X}
    \toprule
    \textbf{Observed mismatch} & \textbf{Why bytes cannot be reused} &
    \textbf{Response}\\
    \midrule
    Different source epoch & The same logical row key can now name different
    sequence bytes & Reject generation; rebuild only when policy permits.\\
    Same \(k,w\), different hash scheme & Query and target keys inhabit
    different hash languages & Reject; no conversion can recover omitted
    minimizers.\\
    Unsupported position width & A decoded occurrence can truncate or wrap &
    Reject before allocation or device upload.\\
    Current GUC names a different provider & The generation's data plane does
    not move when configuration changes & Use the provider stored in its
    descriptor.\\
    Checksum or section-bound failure & Payload is corrupt, misaddressed, or
    incomplete & Reject the selected extent/generation; never continue with
    partial postings.\\
    Strict persisted requirement, no generation & Recompute would change the
    experiment and operational contract & Raise an error rather than falling
    back.\\
    \bottomrule
  \end{tabularx}
\end{table}

This strictness also protects accelerator dispatch. A GPU-resident view is
derived only after host-side format and compatibility validation. The Vulkan
worker receives plain spans and immutable descriptors, not PostgreSQL tuple
objects. If a device view is cached, its key includes the persisted assembly
identity or the exact in-memory shape. A superficially similar generation
cannot alias the old device allocation.

\section{Selective cold loading and memory accounting}
\label{sec:index-cold-loading}

\subsection{The useful numerator is admitted memory}

Reporting the serialized index size is not enough to establish that a database
query is bounded. During a cold read the backend can simultaneously own a
directory, control rows, postings, decoded occurrence arrays, a transient hash
lookup, query-occurrence tables, I/O slabs, and accelerator staging. Some
allocations are visible to PostgreSQL memory contexts, some belong to C++
containers, and some are device memory. A correct budget must cover the
high-water overlap, not just the largest final vector.

BioSeqDB uses explicit memory permits for the significant owners. Format~6
reports separate evidence for directory, control, and posting admission.
Readers authorize an extent before allocating its destination. Merge and
construction code similarly accounts for resident sorted runs and spills the
next run when a permit cannot be obtained. Accelerator-resident objects own
their permit for the complete lifetime of the device allocation rather than
returning it immediately after upload.

The DOC0381 diagnosis illustrates the difference between selective logical
content and inefficient physical access. Its workload selected 17,166 extents
and approximately 251.2 million controls from a generation containing
14,171,377 query hashes. After the occurrence filter it retained 47,948,924
positions and reported 964,465,532 bytes of admitted/owned memory. These
figures show that the selective algorithm was already avoiding a full decoded
position set. Nevertheless, the relation provider performed hundreds of
thousands of small index and heap operations and repeated whole-generation
accounting. Logical selectivity did not compensate for an unsuitable physical
read path.

The corrective design therefore attacks both axes:

\begin{enumerate}
  \item directory and control structures make the authorized postings set
  smaller;
  \item accounting walks each extent a constant number of times instead of
  rescanning the whole set for every batch;
  \item adjacent physical ranges are coalesced into bounded slabs;
  \item checksums are computed at the same granularity as authorization; and
  \item the optional mapped MPF image removes tuple and heap-fetch overhead
  where a complete directly consumable image is the better representation.
\end{enumerate}

\subsection{Backend-local residency caches}
\label{sec:index-residency}

Cold persistence does not imply that every query must reread every selected
chunk. The PostgreSQL backend maintains two bounded LRU caches:

\begin{itemize}
  \item a chunk cache of at most 256 entries and 2 GiB; and
  \item an assembled-index cache of at most 32 entries and 4 GiB.
\end{itemize}

A chunk key includes the full relation and index identity, bucket, shard,
query key, sub-key, and chunk number. An assembled key additionally includes
backend/format shape, scheme, query signature, and clip interval. Values hold
memory permits for their resident bytes. The key breadth is intentional: a
cache entry must not survive a source generation or algorithm change merely
because a short SQL name remains equal.

Transaction, subtransaction, and relcache callbacks maintain coherence. A
write recorded in an aborted transaction is evicted; commit marks eligible
entries stable. Relation invalidation and explicit index-identity invalidation
remove affected entries. This is a backend-local acceleration, not a
cross-backend consistency mechanism: PostgreSQL snapshots and generation
descriptors remain authoritative.

GPU residency follows the same idea at a different boundary. The Phase-B
implementation assembled a GPU view from existing serialized buckets and
kept it in an LRU keyed by the in-memory shape or persisted assembly identity.
That first implementation was 5--8x slower on small cases and 1.8x slower on
its XXL case, demonstrating that residency alone does not erase conversion and
launch overhead. Phase C replaced a problematic open-addressed device table
that could trigger device loss with sorted dense buckets and binary search.
On its later XXL measurement the GPU route took 7.558 s versus 9.501 s for
SIMD, a 1.26x improvement, with 93,309,652 target positions and no fallback.
Those numbers characterize an accelerator front-end, not index persistence,
but they explain why persistent metadata includes enough shape identity to
make safe direct residency possible.

\section{Runtime selection}
\label{sec:index-runtime-selection}

At execution time the scheduler does not simply ask “does a table named
\texttt{mpf} exist?” It resolves an effective index policy for each sequence
column and combines it with the query's read requirement, algorithm
compatibility, source identity, memory admission, and backend availability.
The main decision shape is:

\begin{algorithm}[ht]
  \caption{Persisted-index selection in the alignment dispatcher}
  \label{alg:index-runtime-selection}
  \small
  \begin{algorithmic}[1]
    \Require source identity \(I\), query sketch \(Q\), active index policy \(P\)
    \State resolve the active read name for the source column
    \State inspect a visible persisted generation only if stored reads are enabled
    \If{a generation is present}
      \State validate source epoch, algorithm digest, format, lengths and checksums
      \If{the active name is MPF}
        \State assemble or map the fragment-native lookup under memory permits
      \ElsIf{the active name is MPS}
        \State load the directory, selected buckets/controls, and authorized postings
        \State build the query-specific clipped lookup
      \EndIf
      \State optionally obtain a compatible cached CPU or accelerator view
    \ElsIf{\(P\) requires a compatible persisted index}
      \State fail closed
    \Else
      \State sketch the target and build a transient index
      \State optionally schedule identity-only persistence when all gates permit it
    \EndIf
    \State execute seed lookup; retain the selected generation/view until the batch ends
  \end{algorithmic}
\end{algorithm}

An absent index is thus not always an error, and a present index is not always
usable. This distinction is essential to benchmarks. The \texttt{no\_index}
case should prove that absence was required; the \texttt{just\_index} case
should prove that a compatible persisted generation was used and that
construction was outside the measured region. A fast row without these
attestations can be a recomputation result, a warm assembled cache, or a
different algorithm.

The default active read name is MPF. Stored MPF is assembled from fragment
images or consumed through the supported direct provider. Stored MPS uses its
directory to form a query-specific index. If policy allows recomputation, both
routes can fall back to a transient target index. Deferred persistence scans
the enabled preference list in MPF-then-MPS order; MPF wins if both are
eligible. This deterministic priority prevents two background products from
being presented as one implicit “best index.”

\section{Comparison with minimap2's \texttt{.mmi} index}
\label{sec:index-minimap2-comparison}

The bundled minimap2 source provides a useful algorithmic baseline because the
BioSeqDB seed lookup derives from the same mapping family. Its
\texttt{mm\_idx\_t} records bucket bits \(b\), window \(w\), \(k\), flags,
sequence metadata, an optional packed reference \(S\), and an array \(B\) of
buckets. A bucket contains an array of minimizer/position pairs during
construction, an open-addressing hash table, and a positions array.

Lookup partitions a minimizer hash with the low \(b\) bits selecting the
bucket. The remaining high key is found in that bucket's hash table. A low
marker bit in the stored key distinguishes a singleton: its one position is
stored directly in the hash value. A repeated minimizer stores a packed
\((offset,count)\) value whose occurrences occupy the bucket's \texttt{p}
array. Each occurrence encodes reference identifier, position, and strand.
Construction radix-sorts minimizer/position pairs and sorts the multi-hit
position lists.

The \texttt{.mmi} writer serializes a magic and five 32-bit parameters,
sequence names and lengths, every bucket's position array and hash entries,
and optionally the packed reference. The loader in the bundled implementation
uses \texttt{fread}, allocates all buckets and position arrays, and reconstructs
the complete in-memory index; it is not a selective mmap reader. Index
parameters stored in the file can override command-line values, with a warning.
This format is compact and practical for a command-line mapper whose unit of
reuse is a reference file.

\begin{table}[ht]
  \centering
  \scriptsize
  \caption{Persistent-index responsibilities in bundled minimap2 and BioSeqDB.
  The comparison concerns architecture, not a claim that one wire format is a
  drop-in replacement for the other.}
  \label{tab:mmi-bioseq-index}
  \begin{tabularx}{\linewidth}{p{.19\linewidth}p{.34\linewidth}X}
    \toprule
    \textbf{Concern} & \textbf{minimap2 \texttt{.mmi}} &
    \textbf{BioSeqDB}\\
    \midrule
    Ownership & Standalone file named by the caller &
    Derived object bound to a PostgreSQL physical relation, column, source
    epoch, and generation descriptor.\\
    Logical shape & Bucket-local hash table; singleton inline, repeated
    occurrences in \texttt{p} &
    MPF fragment images or MPS hash-range directory; generation Formats 5/6
    provide specified controls and positions streams.\\
    Load granularity & Reconstruct complete bucket arrays and optional
    reference with \texttt{fread} &
    Full MPF image when appropriate, or query-selected MPS/Format-6 controls
    and postings; bounded admission and backend LRU reuse.\\
    Update & User rebuilds/replaces a file &
    Manual or insert policy; validate unpublished generation, atomic head
    replacement, snapshot-safe predecessor retirement.\\
    Compatibility & File parameters and flags; mapper-specific load contract &
    Alphabet, widths, \(k,w\), scheme/version, producer, occurrence policy,
    source epoch, and optional accelerator ABI digest.\\
    Transactions & Outside database transaction semantics &
    Descriptor publication, abort cleanup, registered snapshots,
    \texttt{ResourceOwner} mappings, and visibility-delayed reaping.\\
    Physical providers & Ordinary file &
    Default unlogged PostgreSQL relation; optional PostgreSQL-governed immutable
    mapped file for supported MPF generations.\\
    SQL batch sharing & One process consumes one loaded index &
    Query batch shares occurrence analysis and assembled views; planner and
    custom scan choose the route.\\
    \bottomrule
  \end{tabularx}
\end{table}

BioSeqDB's legacy MPF layout resembles minimap2 at the conceptual level: a
hash maps to an offset/count and a separate array stores positions. The
important differences are not cosmetic serialization choices. BioSeqDB's
position has a global sequence coordinate rather than minimap2's packed
reference identifier, and the index may begin fragmented because the database
stores long sequences in independently addressable pieces. It must also
authenticate a minimizer scheme shared with Chapter~5's SIMD and Vulkan
implementations. A minimap2 index assumes its mapper controls the surrounding
process; BioSeqDB must make the object safe for SQL transactions, lazy source
reads, memory admission, and multiple execution backends.

\section{Design evolution}
\label{sec:index-design-evolution}

The implemented architecture is easier to understand as a sequence of
constraints discovered by measurement. The chronology in
Table~\ref{tab:index-evolution} does not imply that every mechanism remains
enabled. It records what each stage taught the next one.

\begin{table}[p]
  \centering
  \scriptsize
  \caption{Evolution of BioSeqDB persistent minimizer storage. Newer
  documents and current source supersede older status claims.}
  \label{tab:index-evolution}
  \begin{tabularx}{\textwidth}{p{.105\textwidth}p{.22\textwidth}
                               p{.28\textwidth}X}
    \toprule
    \textbf{Evidence} & \textbf{Design at that point} &
    \textbf{Observed limitation or result} & \textbf{Retained consequence}\\
    \midrule
    DOC0279, DOC0282 &
    Typed registry; MPF migration; MPS introduced &
    One generic ``index blob'' API could not express fragment, sequence, and
    table granularity safely &
    Compile-time algorithm/layout registry and distinct SQL index names.\\
    DOC0288 &
    PostgreSQL ingestion followed by automatic index update &
    Early persisted reuse was competitive on several 2/4 Gbp cases, but the
    old large build could consume 80--93 GB of temporary space &
    End-to-end correctness is necessary; build resource evidence is a separate
    acceptance condition.\\
    DOC0298 &
    Sequence-wide hash directory and query-selected MPS buckets &
    Building and merging all minimizers could retain excessive memory; partial
    output needed a completeness rule &
    Reserved sub-key namespace, fixed directory records, and directory-last
    publication.\\
    DOC0299--DOC0300 &
    Memory-bounded sorted runs and PostgreSQL \texttt{BufFile} spill &
    First measured builder produced 440,117 temporary files and six merge
    passes despite bounded run memory &
    Bounded fan-in, coarser physical units, spillable heavy-hash accumulator,
    and explicit run residency.\\
    DOC0301 &
    Durable, lazy, adaptive, and query-filtered variants &
    Durable A6 completed correctly but took 303.399 s wall; transient A16 took
    about 26.3 s by retaining 0.8078 percent and persisted nothing &
    Persistence attestation is mandatory; query selectivity was retained as a
    design goal, not confused with a completed durable build.\\
    DOC0307--DOC0311 &
    GPU-resident derived view, then dense sorted GPU buckets &
    First residency path was slower; open-addressing caused device loss;
    sorted binary search later beat SIMD on its XXL case &
    Persistent compatibility includes shape/ABI; device representation is
    derived only after host validation and owns a residency permit.\\
    DOC0371, DOC0373 &
    Immutable PostgreSQL generations and shape-aware compact layouts &
    Replace-in-place and generic blob deserialization could not give
    snapshot-safe publication or bounded selective cold reads &
    Generation head, manifests, source epoch, compatibility digest, selective
    controls/postings, unlogged one-page payload rows.\\
    DOC0381 &
    Instrumented cold PostgreSQL relation reader &
    103.840 s index/map region, quadratic accounting, duplicate CRC, and about
    562,000 heap fetches &
    Linear accounting, coalesced ranges, address-bound CRC; reconsider physical
    provider rather than optimizing only logical selectivity.\\
    DOC0383--DOC0384 &
    Fragment-native MPF plus PostgreSQL-governed mapped generation &
    Implemented six-run P50 5.368746 s versus minimap2 16.487527 s, narrowly
    satisfying a 3x gate under documented caveats &
    Default PostgreSQL provider plus opt-in immutable mapped MPF provider,
    transactional descriptor, snapshot-safe retirement, bounded reaper.\\
    DOC0385 &
    Per-query occurrence-merge prototype &
    Slower and higher RSS on every tested workload &
    Keep compact global query-occurrence radix for a batch.\\
    \bottomrule
  \end{tabularx}
\end{table}

\subsection{What the current solution actually selects}

The phrase ``final index format'' would be misleading because the current
system deliberately preserves more than one verified route. The default
policy is nevertheless concrete:

\begin{enumerate}
  \item MPF is the default active minimizer index and is enabled for
  update-on-insert. It matches the fragment-native sequence store and avoids a
  global MPS build.
  \item The default generation data provider is an unlogged PostgreSQL
  relation. The immutable mapped-file provider is implemented and opt-in for
  the supported MPF path.
  \item Format~5 is the generally applicable shape-aware generation format.
  The specialized MPS Format-6 hash-major path is available for its registered
  dna4 \(k=15,w=10\) profile, but the grouped persisted-generation setting is
  disabled by default.
  \item MPS remains a manual index option when query-selective cold loading is
  worth its construction and merge cost.
  \item A strict persisted-read request fails closed. An ordinary request can
  recompute, and may schedule an identity-only persistence job when the
  explicit asynchronous gates allow it.
\end{enumerate}

This is the point at which the BioSeqDB design differs most clearly from its
original minimap2 inspiration. Minimap2 chooses one self-contained reference
index file whose internal buckets are rebuilt into process memory. BioSeqDB
chooses an index \emph{policy} tied to relational identity. It can preserve
fragment-native construction, publish immutable database generations,
authorize only query-relevant postings, account memory across a SQL batch, and
derive a device view. The cost is a larger compatibility and lifecycle surface.

\section{Performance evidence}
\label{sec:index-performance}

\subsection{Early evidence: persistence could help, but did not solve building}

DOC0288 measured a PostgreSQL E2E workflow after the loader's flush/finish
states triggered \texttt{bioseq.update\_index} through SPI and reset the
backend cache. Table~\ref{tab:index-doc0288} reproduces its
\texttt{just\_index} results. They are historical results from that
implementation, not current-provider claims.

\begin{table}[ht]
  \centering
  \small
  \caption{DOC0288 historical \texttt{just\_index} results. Speedup is
  minimap2 time divided by BioSeqDB time. Values are mean
  \(\pm\) reported variation.}
  \label{tab:index-doc0288}
  \begin{tabular}{llrrl}
    \toprule
    \textbf{Target} & \textbf{Alphabet} &
    \textbf{BioSeqDB (s)} & \textbf{minimap2 (s)} & \textbf{Ratio}\\
    \midrule
    2 Gbp & dna4 & \(14.29\pm0.22\) & \(21.53\pm0.30\) & 1.51x\\
    2 Gbp & dna5 & \(15.97\pm0.51\) & \(18.66\pm1.05\) & 1.17x\\
    4 Gbp & dna4 & \(23.20\pm0.65\) & \(32.22\pm0.89\) & 1.39x\\
    4 Gbp & dna5 & \(25.65\pm0.94\) & \(25.48\pm0.55\) & parity\\
    \bottomrule
  \end{tabular}
\end{table}

All checked hits were correct and reported target offsets differed by 2--8 bp.
The same report's \texttt{no\_index} rows were mixed: BioSeqDB was slower at
2 Gbp, but at 4 Gbp it took 171.34 s versus 276.41 s for dna4 and 190.40 s
versus 227.36 s for dna5. The report observed approximately 15 GB of swap for
minimap2 versus about 2.6 GB for BioSeqDB on the 4 Gbp no-index case. These
numbers show the interaction between construction and memory pressure; they do
not isolate a wire format.

The apparent success of the \texttt{just\_index} rows coexisted with the
80--93 GB temporary peaks described earlier. This is why the later work did not
stop at “persisted queries are fast.” It pursued bounded construction,
query-selective reading, immutable publication, and finally a physical provider
that did not turn a multi-gigabyte image into hundreds of thousands of heap
fetches.

\subsection{Intermediate results must retain their persistence label}

\begin{table}[ht]
  \centering
  \scriptsize
  \caption{Selected DOC0300--DOC0301 MPS experiments. The final column prevents
  fast transient experiments from being presented as durable index results.}
  \label{tab:index-mps-intermediate}
  \begin{tabularx}{\linewidth}{p{.12\linewidth}p{.18\linewidth}
                               p{.16\linewidth}p{.16\linewidth}X}
    \toprule
    \textbf{Run} & \textbf{Intent} & \textbf{Wall / region} &
    \textbf{Persistence} & \textbf{Interpretation}\\
    \midrule
    DOC0300 E1f &
    First bounded \texttt{BufFile} build &
    456.715 / 455.452 s &
    durable build path &
    Bounded resident runs, but 33.148 GiB temp accounting and 440,117 files
    made the implementation non-viable.\\
    DOC0301 A6 &
    Valid foreground MPS plus wait &
    303.399 / 91.316 s &
    650 persisted rows &
    Correct 2 bp offset; about 210 s was outside the foreground region while
    persistence completed.\\
    DOC0301 A12 &
    transient no-persist recompute &
    about 89 s &
    none &
    Useful recomputation baseline; not a cold persisted read.\\
    DOC0301 A14 &
    adaptive waves &
    33.797 s wall &
    none &
    Faster scheduling/build shape, but no durable artifact.\\
    DOC0301 A16 &
    query-filtered transient target index &
    26.312--26.454 / 25.071--25.113 s &
    none; audit rows \(=0\) &
    Kept 0.8078 percent of target minimizers and returned the correct 2 bp
    offset. Demonstrates selectivity only.\\
    \bottomrule
  \end{tabularx}
\end{table}

The distinction is central to the thesis's narrative. The fastest intermediate
solution deliberately stopped solving storage. It proved that the query did
not need most target minimizers. The final storage work then had to preserve
that advantage without making construction, integrity verification, or tuple
access dominate.

\subsection{Current mapped-provider evidence}

DOC0384 reports six retained BioSeqDB \texttt{just\_index} samples:
\[
  6.389712833,\ 5.401551458,\ 5.343038250,\ 5.368746000,\
  5.610243750,\ 5.321436000\ \mathrm{s}.
\]
Their median is 5.368746 s. The report's minimap2 median is 16.487527 s, giving
a 3.071020x ratio. Figure~\ref{fig:index-doc0384-bars} displays the medians on
one scale.

\begin{figure}[ht]
  \centering
  \begin{tikzpicture}[x=.60cm,y=.78cm,font=\small,>=Latex]
    \draw[->] (0,0)--(18.3,0) node[right]{seconds};
    \foreach \x in {0,2,4,6,8,10,12,14,16,18} {
      \draw (\x,.08)--(\x,-.08);
      \node[below,font=\scriptsize] at (\x,-.08) {\x};
    }
    \draw[fill=blue!45] (0,.55) rectangle (5.368746,1.25);
    \node[anchor=east] at (-.2,.90) {BioSeqDB};
    \node[anchor=west,font=\scriptsize] at (5.50,.90) {5.369};
    \draw[very thick] (5.321436,1.33)--(6.389713,1.33);
    \draw (5.321436,1.24)--(5.321436,1.42);
    \draw (6.389713,1.24)--(6.389713,1.42);
    \node[anchor=west,font=\scriptsize] at (6.52,1.33)
      {observed BioSeqDB range};

    \draw[fill=orange!60] (0,2.05) rectangle (16.487527,2.75);
    \node[anchor=east] at (-.2,2.40) {minimap2};
    \node[anchor=west,font=\scriptsize] at (16.62,2.40) {16.488};
  \end{tikzpicture}
  \caption{DOC0384 six-run median comparison for the implemented mapped MPF
  path. The thin line shows the range of the six retained BioSeqDB samples; the
  report supplies a minimap2 median rather than a paired range.}
  \label{fig:index-doc0384-bars}
\end{figure}

Correctness was exact for all 16 checked cases in match position, strand, and
CIGAR. The speed result must retain its caveats. The median was only
0.127096 s below the time corresponding to the report's 3x acceptance limit.
The run used a manually assembled, dirty worktree, process-cold operating-system
cache state was not fully specified, and not every individual run passed the
3x boundary. A later fresh E2E row in the same report measured
6.080941750 s versus 43.118770584 s (7.0908x in that same run), but it also
recorded a wrapper/configuration artifact. It is corroboration that the new
path is viable, not a statistically interchangeable replacement for the
six-run distribution.

These qualifications are not incidental. DOC0381 had shown a 6.220x outer
wall-clock disadvantage for the prior cold relation path; DOC0384 shows that a
different physical data plane can reverse the result. Neither number is a
universal property of “indexes in PostgreSQL.” They measure particular source
states, provider paths, and attested execution regions.

\section{Operational consequences and limitations}
\label{sec:index-operational}

Persistent derived state introduces administration that a benchmark process
does not face:

\begin{itemize}
  \item An unlogged relation can be reset after a crash and must be rebuilt
  from the durable source. The system must not advertise its pre-crash
  descriptor as reusable.
  \item The mapped provider needs backup, restore, replication, and tablespace
  procedures that preserve the private artifact together with its descriptor,
  or else deliberately discard and rebuild it.
  \item Source updates need a new epoch and generation. Incremental mutation of
  an image in place would defeat snapshot safety and compatibility checks.
  \item The mapped provider currently solves the supported MPF path; MPS
  asynchronous persistence and specialized Format~6 have narrower gates.
  A benchmark of one route should not be generalized to every index policy.
  \item Backend and device caches reduce warm cost but cannot be used as
  evidence for cold performance unless their state is measured and reported.
  \item Reaping is intentionally bounded and may take several passes. This
  prevents cleanup from becoming an unbounded backend pause, at the price of
  temporary retired generations.
\end{itemize}

There is also a semantic limitation inherited from deriving a minimizer index.
It accelerates seed lookup; it does not turn BioSeqDB into a byte-identical
minimap2 execution. Chapter~5 documents minimizer-scheme differences,
Chapter~6 documents chaining differences, and Chapter~7 documents extension
and CIGAR boundaries. The index compatibility contract ensures that BioSeqDB's
own query and target paths agree. It does not claim that any BioSeqDB
generation can be handed to an arbitrary minimap2 release as a \texttt{.mmi}
file.

\section{Conclusions}
\label{sec:index-storage-conclusions}

BioSeqDB's persistent indexing work began with a conventional idea: serialize
the target's minimizer hash table so the next query can skip rebuilding it.
The database setting exposed the missing parts of that idea. A useful index
must identify the source and algorithm that produced it, publish atomically,
coexist with readers of an older version, survive errors without appearing
complete, load under a memory bound, select only query-relevant postings when
possible, and convert safely into CPU and GPU views.

The final architecture combines several retained decisions:

\begin{enumerate}
  \item typed MPF, MPS, and metadata indexes rather than one opaque blob API;
  \item fragment-native MPF as the default, with MPS available for
  sequence-wide selective loading;
  \item immutable generation descriptors and fail-closed compatibility;
  \item Format~5 as the general shape-aware format and a gated hash-major
  Format~6 route;
  \item a default unlogged PostgreSQL-relation provider and an implemented
  opt-in mapped MPF provider whose lifecycle is governed by PostgreSQL;
  \item explicit memory permits, bounded construction spill, and backend/device
  residency caches; and
  \item deterministic runtime policy and identity-only background persistence.
\end{enumerate}

The intermediate results explain why each piece exists. A global build could
be fast yet consume tens of gigabytes of temporary storage. A bounded builder
could complete yet create hundreds of thousands of files. A query-filtered
transient index could be fast yet solve no persistence problem. A logically
selective PostgreSQL representation could still lose to hundreds of thousands
of heap fetches. The current solution is therefore not simply “store the
minimap2 index in a table.” It separates logical organization, physical
provider, compatibility, publication, and runtime view so each can be measured
and changed without weakening the others.
